\uuid{tIqC}
\exo7id{136}
\titre{exo7 136}
\auteur{cousquer}
\organisation{exo7}
\datecreate{2003-10-01}
\isIndication{false}
\isCorrection{true}
\chapitre{Logique, ensemble, raisonnement}
\sousChapitre{Ensemble}
\cpge{mpsi-mp2i.raisonnement-vocabulaire-ensembliste.b, pcsi.raisonnement-vocabulaire-ensembliste.b, ptsi.raisonnement-vocabulaire-ensembliste.b, mpsi-mp2i.nombres-reels-suites.b, pcsi.nombres-reels-suites.a, ptsi.nombres-reels-suites.a, bcpst1.nombres-reels}
\module{Algèbre}
\niveau{L1}
\difficulte{}

\contenu{
\question{
Montrer que chacun des ensembles suivants est un intervalle,
éventuellement vide ou réduit à un point
$$I_1=\bigcap_{n=1}^{+\infty}\left[3,3+\frac{1}{n^2}\right[\;\mbox{ et
}\; I_2=\bigcap_{n=1}^{+\infty}\left]-2-\frac{1}{n},4+n^2\right].$$
}
\reponse{
$I_1=3\;\mbox{ et
}\; I_2=\left[-2,5\right].$
}
}
